\documentclass[11pt]{article}
\usepackage[a4paper,margin=25mm]{geometry}
\usepackage{amsmath,amssymb,amsthm,booktabs}
\newtheorem{theorem}{Theorem}\newtheorem{lemma}[theorem]{Lemma}
\title{Restricting a normal quotient search to divisors of its target order}
\author{}\date{}
\begin{document}\maketitle
We use the previously reviewed finite-target Coxeter enumeration and the
LINS normal-subgroup algorithm. Their completeness and the supplied finite
group tables remain computational premises. This note justifies a further
restriction when the desired quotient has a prescribed order. It then
applies that restriction to the last square diagram in Table 3 of
Fernandes--Leemans--Weiss, \emph{An Exploration of Locally Spherical
Regular Hypertopes}, fixed 16-page edition, physical page 14.

\begin{lemma}\label{divisor}
Let $N\lhd W$ have finite index $m$. Every subgroup $H$ containing $N$
has index dividing $m$. In particular, every normal overgroup of $N$ is
retained by the condition $[W:H]\mid m$.
\end{lemma}
\begin{proof}
The index formula gives $m=[W:H][H:N]$.
\end{proof}

\paragraph{The search and its dependencies.}
The acquired search first supplies the required normal kernels of finite
target epimorphisms. Its later operations descend through normal
$p$-group sections and take intersections of already known normal
subgroups. The finite-target list, module calculations and kernel
construction are unchanged. We apply the following restrictions:
\begin{enumerate}
\item A finite epimorphic target $Q$ is used only if $|Q|\mid m$.
\item A node $H$ is expanded only if $[W:H]\mid m$.
\item A $p$-section below $H$ is considered only if $p[W:H]\mid m$;
      a proposed elementary section is retained only if its resulting
      index $[W:K]$ divides $m$.
\item An intersection is retained only if its computed index divides $m$.
\end{enumerate}
The ordinary bound remains $m$. No early termination is introduced.
In particular, ``divides $m$'' is not replaced by ``equals $m$'' at an
intermediate step.

To justify completeness, fix a desired $N$ and follow the acquired
algorithm's completeness argument inside the finite normal-subgroup
interval $[N,W]$. Every epimorphic target used in this interval is a
quotient of $W/N$, so its order divides $m$. Every intermediate subgroup
in a descent contains $N$. Its index and each intervening elementary
section therefore satisfy the tests above. If an intersection produces
a subgroup containing $N$, both operands also contain $N$, so neither
operand nor their intersection is removed. Lemma~\ref{divisor} thus
preserves every dependency needed to obtain $N$.

There is also an implementation issue concerning the routine which
decides whether a $p$-section is already obtainable by intersections.
For any retained $H$, every normal supergroup of $H$ has index dividing
$[W:H]$, hence dividing $m$. Its complete supergroup interval is retained.
Consequently the minimal-supergroup relations used by that routine have
not changed. The same applies to the common-supergroup computation for
an intersection. The existing correctness of these optimizations is
therefore inherited on the retained interval. It is not inferred from
small tests alone.

\newpage
\paragraph{Exact implementation and controls.}
The program uses LINS's documented option hooks for branch cuts,
$p$-quotients, proposed modules and intersections. The initial finite-target
routine is the acquired Coxeter routine with one additional order test.
The exact presentation is checked there; its involution choices include
the identity and proper divisors of bond orders, because a quotient on
the way to the desired group may collapse a generator or a bond.
All final marked quotients are tested separately for exact orders.
The program requires a positive target within the installed table bound,
and requires the search bound to equal that same target.

For $A_3,B_3,C_2^3,\widetilde A_2,H_3$, the controls compare the restricted
search with the full default search. There are 26 target choices and
132 retained normal kernels. Each comparison checks the actual two-way
subgroup containment, not merely its index. The $H_3$ cases exercise the
nonabelian finite targets. The elementary abelian cases retain many
distinct kernels of the same index. A deliberately incorrect restriction
to intermediate indices equal to the final target loses kernels, confirming
why normal ancestors must be preserved. Invalid target values and a
mismatched search bound must be rejected before a search can be credited.

\paragraph{The square with labels $3,5,3,5$.}
Put
\[
 W=\langle s_1,s_2,s_3,s_4\mid s_i^2=1,\
 (s_1s_2)^3=(s_2s_3)^5=(s_3s_4)^3=(s_4s_1)^5=1,\
 (s_1s_3)^2=(s_2s_4)^2=1\rangle.
\]
Every finite regular hypertope with this diagram gives a quotient of
this marked presentation. Its four distinguished generators must satisfy
\[
 H_I\cap H_J=H_{I\cap J},\qquad
 H_I=\langle t_i:i\in I\rangle.
\]
This necessary condition applies whether a finite-diagram residue is
required to be a full spherical Coxeter geometry or a projective quotient.

For target $m=9720$, the retained normal-kernel indices are
\[
 1,\quad 2,\quad 60,\quad 120,\quad 4860,\quad 9720,
\]
with one kernel at each index. The unique kernel of index 9720 gives the
27-point marked quotient bound in the literal certificate. It has the
exact required pair orders and four maximal parabolics of order 120.
Nevertheless it fails the intersection property. With indices starting
at 1, the literal closure gives
\[
 |H_{\{1,2,3\}}\cap H_{\{1,2,4\}}|=12,
 \qquad |H_{\{1,2\}}|=6.
\]
There is a shorter obstruction in the same tuple. For each omitted $j$,
form $z_j=(\prod_{i\ne j}t_i)^5$, with the remaining indices increasing.
All four words give the same nonidentity involution $z$. Thus $z$ belongs
to every maximal parabolic, whereas repeated application of the
intersection property would make their total intersection trivial.
The equality and nonidentity of these four words are verified directly
on the 27 points; full group closure is unnecessary for this obstruction.

\begin{theorem}
There is no regular hypertope of square type $(3,5,3,5)$ whose
type-preserving group has order 9720, conditional on the stated normal
quotient enumeration and its retained classification premises.
\end{theorem}
\begin{proof}
Any such hypertope supplies an index-9720 normal kernel of $W$. The
complete restricted enumeration contains exactly one. Its marked quotient
fails the necessary intersection property as shown above. Relabelling
a diagram gives another marking of this same fixed presentation and
does not evade its normal-kernel enumeration.
\end{proof}

\newpage
\paragraph{Attribution, reuse and scope.}
The target order and the displayed group $3^4\rtimes A_5\times C_2$
come from the second $(3,5,3,5)$ entry of the source table. The conclusion
excludes the order itself and hence does not depend on interpreting an
unspecified action in that group name. It is a correction of a published
finite claim, not a solution of an open conjecture.

The normal-subgroup algorithm and its finite group tables are due to
prior work; the installed documentation attributes the underlying
algorithm to Firth's thesis. The finite-target Coxeter traversal, composed
module evaluation, complete marked-intersection test and independent
literal permutation checker are acquired from earlier reviewed components.
The contribution of this note is the explicit justification and use of
the divisor restriction, together with its exact source application.
Divisibility of finite indices is elementary; no historical priority is
claimed for that principle.

The earlier complete search through order 3840 is retained separately.
It contains 25 normal kernels, including eleven of index 3840, and gives
one valid broad marked hypertope of order 1920 with maximal parabolics of
order 60. Its raw finite result and its local interpretation are not
silently incorporated into the theorem here. The general bound 9720
previously timed out; that failed computation is preserved and provides
no completeness premise. The present theorem concerns the exact target
order, not the entire interval below it, and makes no minimum claim.

The certificate checker reruns the bounded restricted computation,
compares the ordered marking with the recorded one by an actual marked
isomorphism, checks the GAP-to-JSON bridge, and applies the independent
integer-permutation closure to all 256 intersections. It separately checks
the four displayed central words and malformed input controls. The proof,
source table and dependency bytes are retained for distinct review.
There is no whole-manuscript credit, new Lean verification, or compression
measurement in this component. Other claims and examples in the source
remain attached to their own recovery obligations.
\end{document}
